A \emph{lattice} is the set of all integer combinations of
linearly independent vectors.
\ifdefined\withModuleSIVP
{Three basic lattice problems are the \emph{shortest vector
problem} (SVP), the \emph{closest vector problem} (CVP), and the
\emph{shortest independent vectors problem} (SIVP).}
\else
Two basic lattice problems are the
\emph{shortest vector problem} (SVP) and the
\emph{closest vector problem} (CVP).
\fi
SVP asks for a shortest nonzero lattice
vector.  CVP asks for a lattice point closest to a given target.
\ifdefined\withModuleSIVP
{In dimension \(n\), SIVP asks for \(n\) linearly independent
lattice vectors whose largest norm is as small as possible.  These}
\else
Both
\fi
 problems are central in the study of lattices and in lattice-based
 cryptography~\cite{MicciancioGoldwasser2002}.

Cryptographic schemes often use lattices with extra algebraic structure.  An
\emph{ideal} in a polynomial quotient ring is closed under
multiplication by every ring element, and a \emph{principal ideal}
is generated by a single element.  In
coefficient coordinates, ring multiplication becomes polynomial convolution.
For example, ideals in \(\Z[X]/(X^D-1)\) give lattices closed under cyclic
rotation.  This structure gives shorter descriptions and faster arithmetic.
It is used in lattice-based cryptography, including the NTRU cryptosystem,
Ring-LWE-based schemes, compact one-way functions, collision-resistant hashing,
and fully homomorphic encryption~\cite{HoffsteinPipherSilverman1998,LyubashevskyPeikertRegev2010,Micciancio2002,PeikertRosen2006,LyubashevskyMicciancio2006,Gentry2009}.  The same structure may
also enable specialized algorithms.

\ifdefined\withModuleSIVP
{Module lattices form an intermediate regime between ordinary
and ideal lattices.  A rank-\(r\) module over the ring of integers of a
degree-\(d\) number field has \(\mathbb Z\)-rank \(rd\).  Rank one gives an
ideal lattice, while taking the base field to be \(\Q\) gives an ordinary
lattice.  Langlois and Stehl\'e introduced Module-SIS and Module-LWE and gave
reductions from worst-case problems on module lattices to these average-case
problems~\cite{LangloisStehle2015}.  Fixed small module rank is now central in
practice, as illustrated by the standardized ML-KEM and ML-DSA
schemes~\cite{NISTFIPS203,NISTFIPS204}.}
\fi

For general lattices, much is known about hardness.  Exact
Euclidean decision-CVP is \(\NP\)-complete~\cite{vanEmdeBoas1981}.  CVP is
also \(\NP\)-hard to approximate within almost-polynomial
factors~\cite{DinurKindlerSafra1998}.  Exact Euclidean SVP is
\(\NP\)-hard under randomized reductions~\cite{Ajtai1998}.  Approximating
Euclidean SVP within some constant factor is also \(\NP\)-hard under randomized
reductions~\cite{Micciancio1998}.  Hair and
Sahai recently proved deterministic \(\NP\)-hardness of exact SVP in the
\(\ell_p\)-norm for every \(p>2\)~\cite{HairSahaiSTOC26}, and
almost-polynomial-factor hardness for every finite \(p\geq2\) under
deterministic subexponential-time reductions~\cite{HairSahaiFOCS26}.
These reductions can
build arbitrary lattice bases and need not preserve the ideal or principal
ideal structure.  Earlier work gave reductions between problems on ideal
lattices, but these reductions do not prove the \(\NP\)-hardness of exact CVP
on principal ideals~\cite{PeikertRosen2007}.  This gap matters because
cyclotomic structure can lead to special algorithms.  Examples include
recovering short generators of some principal ideals~\cite{CramerEtAl2016}
and, under number-theoretic assumptions, finding mildly short vectors in
cyclotomic ideal lattices in quantum polynomial
time~\cite{CramerDucasWesolowski2021}.

\ifdefined\withModuleSIVP
{Classical hardness results are also known for SIVP on general
lattices~\cite{BlomerSeifert1999,AggarwalChung2021}.  For structured
lattices, random self-reducibility was first established for
approximate Ideal-SVP~\cite{deBoerEtAl2020}.  Under ERH, this was later
extended to approximate SIVP on module lattices of any fixed
rank~\cite{deBoerEtAl2025}, yielding a polynomial-time worst-case-to-average-case
reduction over a discretized invariant distribution with a polynomially
bounded loss in the approximation factor.
The exact worst-case complexity of SIVP on explicit fixed-rank module
families is a separate question.}
\fi

\emph{CVP with preprocessing} (CVPP) asks what happens when the lattice is known
before the target.  The offline step may take unlimited time but store only
polynomially many bits, while the online step must answer target-and-threshold
queries in polynomial time.  Micciancio showed that exact
decision-CVPP has no polynomial-time solution on general lattices unless
\(\NP\subseteq\Ppoly\)~\cite{Micciancio2001}.  For approximate CVPP, constant-factor
approximation is hard~\cite{FeigeMicciancio2004}.  Stronger inapproximability
factors are known under quasipolynomial-time hardness
assumptions~\cite{AlekhnovichEtAl2005,KhotPopatVishnoi2012}.  Separately,
Micciancio asked whether CVP
is \(\NP\)-hard on cyclic lattices and whether it remains hard for a fixed
family with one cyclic lattice for each input size~\cite{Micciancio2007}.

\ifdefined\withModuleSIVP
{This leads to our main questions.}
\else
This leads to our main question.
\fi

\begin{quote}
\emph{Is exact Euclidean CVP \(\NP\)-hard for ideal and cyclic lattices, even
for full-rank principal ideals?  Does this hardness still hold for a
uniformly computable family with one ideal per input size, fixed before the
target is known?}
\end{quote}

\ifdefined\withModuleSIVP
{In the module setting, we additionally ask whether exact SIVP
remains hard when the module rank is fixed at two.}
\fi

\subsection{Our results}
\label{sec:intro-results}

\ifdefined\withModuleSIVP
{Our answer to these questions is yes.}
\else
Our answer to both questions is yes.
\fi
Our main result concerns principal ideals in the power-of-two cyclotomic rings
\[
 R_d\mathrel{:=}\Z[y]/(y^d+1),
\]
where \(d\) is a power of two.
\ifdefined\withModuleSIVP
{We write \(K_d\mathrel{:=}\Q[y]/(y^d+1)\) for the fraction
field of \(R_d\).}
\fi

\begin{theorem}[{\normalfont Principal cyclotomic ideals, informal}]
\label{thm:intro-principal-ideal}
For \(d\) ranging over powers of two, exact Euclidean decision-CVP is
\(\NP\)-complete on the coefficient lattices of nonzero principal ideals in
\(R_d\).
\end{theorem}

We use a deterministic polynomial-time many-one reduction
from Exact Cover by 3-Sets (X3C).  The reduction produces a nonzero principal
ideal \((g)\subseteq R_d\), a target \(T\in R_d\), and a squared threshold
\(\Delta\in\mathbb Z_{\geq0}\).  Its closest squared distance satisfies
\[
\begin{aligned}
 {\dist(T,(g))^2}
   &=\Delta &&\text{in a YES instance},\\
 {\dist(T,(g))^2}
   &\geq\Delta+4 &&\text{in a NO instance}.
\end{aligned}
\]
Moreover, the ideal elements at squared distance at most \(\Delta\) from the
target under the coefficient embedding are in bijection with the exact covers
of the X3C instance.

Every nonzero ideal in \(R_d\) has full \(\mathbb Z\)-rank.
By Lemma~\ref{lem:canonical}, the theorem also holds under the
canonical embedding after multiplying all squared distances and thresholds
by \(d\).
This bijection also gives
\(\NP\)-hardness of exact search-CVP under polynomial-time Turing
reductions.

The construction also lifts to full-rank principal ideals in
\(S_D\mathrel{:=}\Z[X]/(X^D-1)\), where
\(D\mathrel{:=}2d\), and proves \(\NP\)-completeness of
exact decision-CVP on their coefficient lattices, which are cyclic
lattices.  Exact search-CVP on these
lattices is \(\NP\)-hard under polynomial-time Turing reductions.

For each X3C universe size \(m\), we also give one uniformly
computable principal cyclotomic ideal and one principal cyclic ideal, both of
dimension \(\Theta(m^7)\).  Only the target and squared
threshold depend on the triple collection.  Exact decision-CVP remains
\(\NP\)-complete on both fixed families.
Unless \(\NP\subseteq\Ppoly\), exact decision-CVPP has no
polynomial-time solution on either family.  Consequently, the polynomial
hierarchy would collapse to \(\Sigma_2^{\mathsf P}\) if such a solution
existed~\cite{KarpLipton1980}.  To our knowledge, these
cyclic results answer Micciancio's two open questions about exact CVP.

\ifdefined\withModuleSIVP
We further transfer the principal cyclotomic CVP instances to exact
module-SIVP on free rank-two modules over \(R_d\).
{For every such module \(\mathcal M\), the coefficient lattice
\(\mathcal L(\mathcal M)\) has dimension \(2d\), and}
\[
 {
 \lambda_2^{K_d}(\mathcal M)
 =\lambda_{2d}\bigl(\mathcal L(\mathcal M)\bigr).
 }
\]
Thus exact decision-module-SIVP coincides with standard exact
decision-SIVP on the underlying coefficient lattice.

\begin{theorem}[{\normalfont Rank-two cyclotomic modules, informal}]
\label{thm:intro-module-sivp}
For \(d\) ranging over powers of two, exact decision-module-SIVP is
\(\NP\)-complete on free rank-two modules over \(R_d\).
\end{theorem}
\fi

\subsection{Technical overview}
\label{sec:intro-overview}

The main challenge is to control every element of the ideal.
In a reduction to a general lattice, we may choose the basis vectors
separately.  For a principal ideal \((g)\subseteq R_d\)
generated by \(g\), however, the natural basis vectors are
the \emph{negacyclic shifts}
\[
 g,yg,\ldots,y^{d-1}g,
\]
and every element of \((g)\) has the form \(gq\) for some
\(q\in R_d\).  Thus we must encode Boolean vectors by suitable choices of
\(q\) and rule out every other \(gq\).

We achieve our goal by reducing from X3C~\cite{GareyJohnson}.
An X3C instance consists of a universe of size \(m\) and a
collection of \(n\) three-element subsets.  It asks whether some subcollection
covers every element exactly once.  We express this condition as an integer
linear system in Boolean variables:
\begin{equation}
\label{eq:intro-x3c-system}
 \mathbf A\boldsymbol\xi=\mathbf b,
 \qquad
 \mathbf A\in\{0,1\}^{m\times n},\quad
 \mathbf b\in\{0,1\}^m,\quad
 \boldsymbol\xi\in\{0,1\}^n.
\end{equation}
Here \(\mathbf A\) is the incidence matrix.  Its entry
\(A_{ij}\) equals one exactly when the \(i\)-th universe element belongs to
the \(j\)-th set.  We take
\(\mathbf b\mathrel{:=}\mathbf 1_m\).  A Boolean solution
\(\boldsymbol\xi\) specifies an
exact cover.  We encode \(\boldsymbol\xi\) as
a polynomial with Boolean coefficients:
\[
 a_{\boldsymbol\xi}(x)
 :=\sum_{k=1}^n\xi_kx^{\alpha_k}.
\]

\paragraph{Encoding the Boolean equations.}
We choose the exponents \(\alpha_1,\ldots,\alpha_n\) so that
all nonzero ordered differences \(\alpha_k-\alpha_j\)
are pairwise distinct.
A classical Sidon
construction provides such exponents of size
\(O(n^2)\)~\cite{ErdosTuran41}.
We choose pairwise well-separated coefficient positions
\(\rho_1,\ldots,\rho_m\) for the \(m\) rows and set
\[
 h(x):=\sum_{i,j:A_{ij}=1}x^{\rho_i-\alpha_j}.
\]
The product
\(x^{\rho_i-\alpha_j}x^{\alpha_k}\) contributes at coefficient position
\[
 \rho_i-\alpha_j+\alpha_k.
\]
When \(k=j\), all contributions for row \(i\) occur at
\(\rho_i\) and sum to \((\mathbf A\boldsymbol\xi)_i\).  When \(k\neq j\),
these products occupy pairwise distinct positions, all separate from the row
positions \(\rho_{i'}\).

To obtain a gap in coefficient \(\ell_2\)-distance, we define the
target polynomial \(t\) by assigning coefficient \(2b_i\) at
\(x^{\rho_i}\) and coefficient one at every position
\(\rho_i-\alpha_j+\alpha_k\) with \(A_{ij}=1\) and \(k\neq j\).  The
positions \(\rho_i\) contribute
\(4\|\mathbf A\boldsymbol\xi-\mathbf b\|_2^2\).  At every other specified
position the difference is \(2\xi_k-1\in\{-1,1\}\).  We denote the number of
these positions by \(\nu\), so their total squared contribution is \(\nu\),
independently of \(\boldsymbol\xi\).  This gives
\[
 \|2h(x)a_{\boldsymbol\xi}(x)-t(x)\|_2^2
 =\nu+4\|\mathbf A\boldsymbol\xi-\mathbf b\|_2^2.
\]
Thus every Boolean vector satisfying
\eqref{eq:intro-x3c-system} has squared distance \(\nu\), whereas every Boolean
vector violating at least one constraint in \eqref{eq:intro-x3c-system} has
squared distance at least \(\nu+4\).

\paragraph{Controlling every ideal element.}
The identity above applies only to the polynomials
\(a_{\boldsymbol\xi}\) encoding Boolean vectors.  To handle every
\(q\in R_d\), we place the construction in the even coordinates of a larger
cyclotomic ring by replacing \(x\) with \(y^2\).  The degree-\(<d\)
representative of \(q\) then decomposes uniquely as
\[
 q(y)=a(y^2)+yb(y^2).
\]
Here \(a\) and \(b\) are the integer polynomials formed by the
even and odd coefficients of \(q\), respectively.  Define
\(u(x):=\sum_{k=1}^n x^{\alpha_k}\).
For a sufficiently large integer \(P\), we set
\[
 \begin{aligned}
 g(y)&:=2h(y^2)+2Py,\\
 T(y)&:=t(y^2)+Py\,u(y^2),\\
 \Delta&:=P^2n+\nu.
 \end{aligned}
\]
Separating \(gq-T\) into its even- and
odd-degree coefficients gives
\[
 gq-T=(2ha+2Pxb-t)(y^2)+y(2hb+2Pa-Pu)(y^2).
\]

The terms involving \(P\) in the two components above have
squared norm
\(4P^2(\|b\|_2^2+\|a-u/2\|_2^2)\).  Over integral
coefficients, this quantity is minimized exactly when \(b=0\) and
\(a=a_{\boldsymbol\xi}\) for some Boolean
\(\boldsymbol\xi\in\{0,1\}^n\).  Every other integral pair increases this
squared norm by at least \(4P^2\).  We choose \(P\) large enough that the
remaining terms cannot offset this increase, making every other \(q\) have
squared distance greater than \(\Delta+4\).

For a Boolean vector \(\boldsymbol\xi\in\{0,1\}^n\), the two components give
\[
 \|g a_{\boldsymbol\xi}(y^2)-T\|_2^2
 =\Delta
  +4\|\mathbf A\boldsymbol\xi-\mathbf b\|_2^2.
\]
Thus the elements of \((g)\) at squared distance at most
\(\Delta\) from \(T\) are exactly those
encoding solutions, and every other element has squared distance at least
\(\Delta+4\).

Full rank follows because \(R_d\) is an integral domain and the generator is
nonzero.  Fourier orthogonality shows that squared distances under
the canonical embedding are \(d\) times those under the coefficient
embedding.

\paragraph{Cyclic and fixed-family extensions.}
For the cyclic case, we convert each constructed principal
cyclotomic ideal into a principal cyclic ideal of twice the dimension, while
multiplying the squared distances of corresponding elements by eight.
For the fixed families, we use the matrix of all
\(\binom m3\) triples.  A mask \(\boldsymbol\mu\) indicates which triples
belong to the input collection, and we replace \(u(x)\) by
\(u_{\boldsymbol\mu}(x):=\sum_k\mu_kx^{\alpha_k}\).  With this replacement,
the odd-coordinate penalty is minimized only if
\(\boldsymbol\xi\leq\boldsymbol\mu\).  Hence absent triples cannot be selected.
The resulting principal cyclotomic and cyclic ideals depend only
on \(m\).  The details appear in \cref{sec:cyclic-cvp} and
\cref{sec:fixed}, respectively.

\ifdefined\withModuleSIVP
\paragraph{The module-SIVP extension.}
We start with an exact-CVP instance \((g,T,\Delta)\) on a principal
cyclotomic ideal.
Here \(g\in R_d\setminus\{0\}\) generates the ideal \((g)\), \(T\in R_d\) is
the target, and \(\Delta\in\mathbb Z_{\geq0}\) is the squared-distance
threshold.  Choose an integer \(Q\) with
\(Q^2>\max\{\|T\|_2^2,\|g\|_2^2\}\) and form
\[
 \mathcal M_{Q,g,T}:=R_d(Q,T)+R_d(0,g)\subseteq K_d^2.
\]
Thus \(\mathcal M_{Q,g,T}\) is a free rank-two
\(R_d\)-module, while its coefficient lattice
\(\mathcal L(\mathcal M_{Q,g,T})\) is a rank-two cyclotomic module lattice.
The submodule \(R_d(0,g)\) spans a one-dimensional
\(K_d\)-subspace, so any two module vectors linearly independent over \(K_d\)
must include a vector outside this subspace.  The choice of \(Q\) forces a
shortest such vector to have the form \(\pm y^j(Q,T)+b(0,g)\) for some
\(b\in R_d\) and \(0\leq j<d\).
Minimizing the norm of its second coordinate then gives
\[
 \bigl(\lambda_2^{K_d}(\mathcal M_{Q,g,T})\bigr)^2
 =Q^2+\dist(T,(g))^2.
\]
Thus the squared-distance gap transfers directly to exact module-SIVP.
\fi


\paragraph{Organization.} \Cref{sec:pre} reviews lattices and
ideal lattices.  \Cref{sec:compiler} proves \(\NP\)-completeness of exact
decision-CVP on principal cyclotomic and principal cyclic ideal lattices.
\Cref{sec:fixed} proves that exact decision-CVPP has no polynomial-time
solution on the resulting fixed principal-ideal families unless
\(\NP\subseteq\Ppoly\).
\ifdefined\withModuleSIVP
\Cref{sec:module-sivp} proves \(\NP\)-completeness of exact
module-SIVP on free rank-two cyclotomic modules.
\fi
\Cref{sec:conclusion} summarizes the results.
